% AUTO-GENERATED by the update_record tool. DO NOT EDIT.
% verify_paper_values regenerates this file from record.json and the
% run outputs and fails on any difference, so manual edits always surface.
\noindent\textbf{H1.} 初等的な総和恒等式（ガウスの和と奇数の和）は、Lean 4 と mathlib の標準公理のみに依存して数学的帰納法により形式証明でき、その検証は実験と同じ事前登録・隔離実行・provenance 照合の流れで行える。
\begin{enumerate}
\item[\textbf{C1}] ガウスの和 2·Σ_{i<n+1} i = n(n+1) は、すべての自然数 n について Lean 4 で sorry 無しに証明され、依存公理は標準 3 公理に収まる。
  \emph{Rationale:} h1 の『初等的な総和恒等式が標準公理のみで形式証明できる』の第一の実例。等差数列の和という最も基本的な恒等式で、帰納法と Finset.sum_range_succ だけで閉じるはず。
  \emph{Verdict:} pending.
\item[\textbf{C2}] 奇数の和 Σ_{i<n} (2i+1) = n² は、すべての自然数 n について Lean 4 で sorry 無しに証明され、依存公理は標準 3 公理に収まる。
  \emph{Rationale:} h1 の第二の実例。c1 と独立な恒等式で、同じ帰納法の型で閉じることが、手法が特定の恒等式に依存しないことの証拠になる。
  \emph{Verdict:} pending.
\end{enumerate}
\noindent\emph{Assumed, so that the claims together imply H1:}
\begin{itemize}
\item c1 と c2 の statement が、それぞれの恒等式の意図した数学的内容を忠実に形式化している（形式化の忠実性は Lean が検証しない）。
\item Lean 4 v4.33.1 と mathlib v4.33.1 のカーネルおよび report ツール（airas-report）が正しく、collectAxioms の返す公理集合が完全である。
\item GitHub Actions 上の実行が、record が宣言した commit の子孫で行われている（provenance 照合が依拠する前提）。
\end{itemize}
